\chapter{Examples}\label{sec:examples}
%import inspect
%import gretl.examples.estimators.ar
%print(inspect.getsource(gretl.examples.estimators.ar.run_example))
%
%import gretl.examples.estimators.ar
%gretl.examples.estimators.ar.run_example()




%%%%%%%%%%%%%%%%%%%%%%%%%%%%%%%%%%%%%%%%%%%%%%%%%%%%%%%%%%%%%%%%%%%%%%%%%%%%%%%%%%
%%%%%%%%%%%%%%%%%%%%%%%%%%%%%%%%%%%%%%%%%%%%%%%%%%%%%%%%%%%%%%%%%%%%%%%%%%%%%%%%%%
%%%%%%%%%%%%%%%%%%%%%%%%%%%%%%%%%%%%%%%%%%%%%%%%%%%%%%%%%%%%%%%%%%%%%%%%%%%%%%%%%%
\section{Examples from \pkg{statsmodels}}
In the \texttt{`examples/statsmodels'} directory, examples
for the \pkg{statsmodels} package, available at
\url{https://www.statsmodels.org/stable/examples},
are provided as \gretlpy calls. It is important to note
that any \python script must be executed from within the
subdirectory where the script is located. This requirement
arises because each example imports functions from
\texttt{helper\_func.py'}, which is located in the
\texttt{`examples/statsmodels'} directory.

Examples available at \url{https://www.statsmodels.org/dev/examples/}
are divided into numbered blocks, with numbers enclosed
in square brackets (e.g., ``[1]:''). The \gretlpy implementation
of \pkg{statsmodels} examples retains this numbering format.
For instance, in the OLS example provided at
\url{https://www.statsmodels.org/dev/examples/notebooks/generated/ols.html},
the ``fit and summary'' output from the first estimation
appears under the number ``[5]:''. This is illustrated in
Figure~\ref{fig:statsmodels_numbering_example_stats}.

\begin{figure}[h!tb]
  \centering
  \includegraphics[width=.8\columnwidth,keepaspectratio]%
    {examples/statsmodels/statsmodels_numbering_example_stats.png}
  \caption{Numbering in \pkg{statsmodels} output: example from \pkg{statsmodels}.}
  \label{fig:statsmodels_numbering_example_stats}
\end{figure}

The corresponding output generated by \gretlpy is shown in
Figure~\ref{fig:statsmodels_numbering_example_gretl4py}.

\begin{figure}[h!tb]
  \centering
  \includegraphics[width=.8\columnwidth,keepaspectratio]%
    {examples/statsmodels/statsmodels_numbering_example_gretl4py.png}
  \caption{Numbering in \pkg{statsmodels} output: example from \gretlpy.}
  \label{fig:statsmodels_numbering_example_gretl4py}
\end{figure}

Currently, we provide examples for:
\begin{itemize}
  \item Linear Regression Models:
    \begin{itemize}
      \item Ordinary Least Squares
        (\url{https://www.statsmodels.org/dev/examples/notebooks/generated/ols.html}),\newline
        file: \texttt{examples/statsmodels/Linear Regression Models/ols.py}.
    \end{itemize}
  \item Time Series Analysis:
    \begin{itemize}
      \item ARMA: Sunspots Data
        (\url{https://www.statsmodels.org/dev/examples/notebooks/generated/tsa_arma_0.html}),\newline
        file: \texttt{examples/statsmodels/Time Series Analysis/arima\_sunspot.py}.
      \item ARMA: Artificial Data
        (\url{https://www.statsmodels.org/dev/examples/notebooks/generated/tsa_arma_1.html}),\newline
        file: \texttt{examples/statsmodels/Time Series Analysis/arima\_artificial.py}.
    \end{itemize}
\end{itemize}
